%! Unit 1, Lessons 1.0–1.7 — ACTUAL Unit Test (Answer Key)
% The scored unit test. Item for item it is the parallel form of ../practice_test_key/: the same
% 15 multiple-choice targets in the same order and the same five free-response questions with the
% same parts and points (MC 15 x 4 = 60, FRQ 5 x 8 = 40, total 100; 20 questions; 5 pages at
% most) --- every context different, so a student who worked the practice test has practised the
% moves and not memorised the answers. NEVER merged into any packet (shared/test_keys.mk drops
% only the practice key).
% NAME ROW: \namedateperiod IS carried --- a unit test is taken in a testing setting.
% NO SKETCH ITEMS: every display is pre-drawn in TikZ and read.
% COMPACT LAYOUT: short answer choices run in two columns (multicols), figures sit beside their
% stems, and each free-response part gets only the lines its answer needs. Keep it to 5 pages.
%
% BLUEPRINT
%   MC 1 grouping (binning)               reading speed in words/min       1.1
%   MC 2 groups of unequal size           two school bands, brass          1.2
%   MC 3 discrete vs continuous           backpack weight to 0.1 lb        1.3
%   MC 4 read a stem plot (median)        15 yoga-class ages               1.3
%   MC 5 interval width                   60 lake depths                   1.3
%   MC 6 shape from a histogram           40 quiz scores                   1.4
%   MC 7 unusual features (dotplot)       20 students' push-ups            1.4
%   MC 8 mean vs median from shape        retirement ages                  1.5
%   MC 9 resistance (a typo)              9 runners' times                 1.5
%   MC10 standard deviation in context    40 water bottles                 1.5
%   MC11 percentile                       a baby's length                  1.5
%   MC12 every section holds a quarter    40 days at a skate park          1.6
%   MC13 parallel boxplots                tomato plants, two fertilizers   1.6
%   MC14 parameter vs statistic           university GPA                   1.7
%   MC15 z-scores across distributions    math vs spelling contests        1.7
%   FRQ1 individuals, variables, SQ       a marina's boats                 1.0, 1.1, 1.3
%   FRQ2 relative frequency comparison    two gyms' favorite classes       1.2
%   FRQ3 histogram, median, describe      50 pumpkin weights               1.3, 1.4
%   FRQ4 five numbers, outlier, center    15 podcast-episode lengths       1.5, 1.6
%   FRQ5 empirical rule, z-scores         coffee-machine fills             1.7
% ANSWERS (MC): 1A 2E 3B 4C 5A 6D 7E 8B 9A 10D 11D 12E 13C 14C 15B
%
% THE DATA (every number checked)
%  MC2  Lincoln 27/90 = 0.30; Monroe 21/60 = 0.35.
%  MC4  23 25 28 30 31 34 34 36 39 42 45 45 48 51 57; median (8th) 36; mean 568/15 = 37.9.
%  MC6  bins of 5 from 20: 1 1 3 6 12 17 (40), skewed left.
%  MC7  20:1 22:2 24:4 26:6 28:4 30:2 44:1 (20).
%  MC9  21 22 24 25 26 27 27 29 31 -> 2.1: Q1 23, med 26, Q3 28 both ways; mean, range move.
%  MC12 12/20/31/36/48.
%  MC13 X 30/38/45/52/60 (IQR 14); Y 35/47/52/57/62 (IQR 10).
%  MC15 Leo (62-50)/8 = 1.5; Kim (45-33)/10 = 1.2.
%  FRQ2 East 105/90/75/30 (300) = 35/30/25/10%; West 72/36/54/18 (180) = 40/20/30/10%.
%  FRQ3 bins of 4 lb from 4: 6 14 12 8 5 3 0 2 (50); cum 6 20 32 40 45 48 48 50.
%  FRQ4 4 22 25 27 28 30 31 32 33 35 36 38 40 41 44: Q1 27, med 32, Q3 38, IQR 11,
%       fences 10.5 and 54.5; sum 466, mean 31.07.
%  FRQ5 N(12, 0.4): 11.6-12.4 68%; <11.2 2.5%; 12.7 -> z 1.75; vs machine 2 N(16, 0.6) 17.2 -> 2.0.
%
% LOCKSTEP: the blank ../../tests/actual_test/ is generated FROM THIS FILE (edit here, regenerate).
% Every ans mark on an MC option prints bare there; every other ans mark becomes a phantom of
% itself; every run of n anslines becomes writelines{n}; work blocks stay byte-identical.
\documentclass[12pt]{article}
\usepackage{apstats-article}
\usepackage{apstats-key}   % pulls in -boxes; do NOT also load -boxes
\usepackage{multicol}
\setlength{\multicolsep}{1pt}

% Answer choices: tight list; wrap in multicols{2} when every choice fits half a line.
\newenvironment{choices}
  {\begin{enumerate}[label=\textbf{(\Alph*)}, leftmargin=2.2em, itemsep=0pt, topsep=1pt]}
  {\end{enumerate}}

\begin{document}

\pageheader{Unit 1 \quad Lessons 1.0--1.7}{Unit Test --- Answer Key}

\namedateperiod
\par\vspace{2pt}
\hfill\small Score:\;\blank{1.3cm}\;/\;100\normalsize
\vspace{0.04in}

\boxguard[6]
\begin{headlinebox}{goldbg}
\small
\textbf{Section I --- Multiple Choice} \hfill \textbf{60 points, 4 each}\\
Notes closed. Circle the single best answer.
\end{headlinebox}

\vspace{0.06in}

\begin{enumerate}[leftmargin=1.9em, itemsep=0.03in]

  \boxguard[8]
  \item A teacher records each student's reading speed in words per minute, then replaces it
        with a label: \emph{Slow} (under 150), \emph{Average} (150 to 250), or \emph{Fast} (over
        250). Which statement is true?
        \begin{choices}
          \item \ans{It is now categorical, and the mean reading speed can no longer be found.}
          \item The new variable is still quantitative, because the labels are based on speed.
          \item The new variable is quantitative, because the labels are in order.
          \item Nothing is lost: the speeds can be recovered from the labels.
          \item The median reading speed can still be found exactly from the labels.
        \end{choices}

  \boxguard[8]
  \item In the Lincoln school band, 27 of 90 members play brass. In the Monroe school band, 21 of
        60 do. Which is best supported?
        \begin{multicols}{2}
        \begin{choices}
          \item Lincoln: 27 is greater than 21
          \item Equally common in both bands
          \item Can't compare: bands differ in size
          \item Lincoln: it has more members
          \item \ans{Monroe: 35\% is greater than 30\%}
        \end{choices}
        \end{multicols}

  \boxguard[6]
  \item Which of these variables is \textbf{continuous}?
        \begin{multicols}{2}
        \begin{choices}
          \item Number of books in a backpack
          \item \ans{Backpack weight, to 0.1 lb}
          \item Number of classes taken
          \item Number of absences this year
          \item Number of pencils in a case
        \end{choices}
        \end{multicols}

  \boxguard[8]
  \item \begin{minipage}[t]{0.62\linewidth}\raggedright
        The stem plot shows the ages (years) of the 15 people in an evening yoga class. What is
        the \textbf{median} age?
        \begin{multicols}{3}
        \begin{choices}
          \item 6 years
          \item 34 years
          \item \ans{36 years}
          \item 37.9 years
          \item 39 years
        \end{choices}
        \end{multicols}
        \end{minipage}\hfill
        \begin{minipage}[t]{0.32\linewidth}
        \centering\small
        \begin{tabular}[t]{@{} r | l @{}}
        2 & 3 5 8 \\
        3 & 0 1 4 4 6 9 \\
        4 & 2 5 5 8 \\
        5 & 1 7 \\
        \end{tabular}\par\vspace{2pt}
        {\footnotesize Key: $3\,|\,6 = 36$ years}
        \end{minipage}

  \boxguard[8]
  \item Two histograms of the \emph{same} 60 lake depths (feet): with 2-foot intervals there are
        two clear peaks; with 10-foot intervals there is one. Which is correct?
        \begin{choices}
          \item \ans{Both are correct: the interval width changes the picture, not the data.}
          \item One of the two histograms must have been drawn incorrectly.
          \item The wider intervals changed some of the depth measurements.
          \item Only the histogram with fewer bars is a true histogram.
          \item The two peaks show that the data set contains two outliers.
        \end{choices}

  \boxguard[9]
  \item \begin{minipage}[t]{0.56\linewidth}\raggedright
        The histogram shows the scores of 40 students on a 50-point quiz. Which best describes
        the \textbf{shape}?
        \begin{choices}
          \item Skewed right: most scores pile up on the right.
          \item Roughly symmetric: there is one peak.
          \item Uniform: every interval has a student.
          \item \ans{Skewed left: a few scores are low.}
          \item Bimodal: the two highest bars are tall.
        \end{choices}
        \end{minipage}\hfill
        \begin{minipage}[t]{0.40\linewidth}
        \centering
        \begin{tikzpicture}[x=0.135cm, y=0.135cm, baseline=(current bounding box.north)]
          \foreach \y in {0,5,10,15} {
            \draw[linegray, very thin] (20,\y) -- (51,\y);
            \node[left, font=\tiny] at (20,\y) {\y};}
          \foreach \l/\c in {20/1, 25/1, 30/3, 35/6, 40/12, 45/17}
            \draw[fill=skymid] (\l,0) rectangle (\l+5,\c);
          \draw[thick] (20,0) -- (52,0);
          \draw[thick] (20,0) -- (20,18);
          \foreach \x in {20,25,...,50} { \node[below, font=\tiny] at (\x,0) {\x}; }
          \node[font=\tiny] at (36,-4.5) {Quiz score (points)};
          \node[rotate=90, font=\tiny] at (14.5,9) {Count};
        \end{tikzpicture}
        \end{minipage}

  \boxguard[9]
  \item \begin{minipage}[t]{0.56\linewidth}\raggedright
        The dotplot shows how many push-ups each of 20 students did in one minute. Which
        statement is best?
        \end{minipage}\hfill
        \begin{minipage}[t]{0.40\linewidth}
        \centering
        \begin{tikzpicture}[x=0.22cm, y=0.2cm, baseline=(current bounding box.north)]
          \draw[thick] (18,0) -- (46,0);
          \foreach \x in {18,22,...,46} {
            \draw (\x,-0.4) -- (\x,0.4); \node[below, font=\tiny] at (\x,-0.3) {\x};}
          \foreach \x/\n in {20/1, 22/2, 24/4, 26/6, 28/4, 30/2, 44/1}
            \foreach \k in {1,...,\n} \fill[navy] (\x,\k*0.9) circle (2pt);
          \node[font=\tiny] at (32,-3) {Push-ups in one minute};
        \end{tikzpicture}
        \end{minipage}
        \begin{choices}
          \item The student with 44 push-ups is an outlier, so that value should be deleted.
          \item There are no unusual features, because most students did 22 to 30 push-ups.
          \item The distribution is bimodal, with peaks at 26 and 44 push-ups.
          \item The gap from 20 to 22 push-ups is the most unusual feature.
          \item \ans{44 push-ups is a possible outlier, set off by a gap; investigate it.}
        \end{choices}

  \boxguard[8]
  \item The ages at which 120 employees of a company retired are strongly \textbf{skewed left}.
        Which is most likely true?
        \begin{multicols}{2}
        \begin{choices}
          \item Mean $>$ median; report the median
          \item \ans{Mean $<$ median; report the median}
          \item Mean $=$ median
          \item Mean $<$ median; report the mean
          \item Can't tell without the data
        \end{choices}
        \end{multicols}

  \boxguard[8]
  \item Nine runners finished a 5-mile race in 21, 22, 24, 25, 26, 27, 27, 29, and 31 minutes. The
        21 is mistakenly entered as 2.1. Which summaries change?
        \begin{multicols}{2}
        \begin{choices}
          \item \ans{Mean and range; not median or IQR}
          \item Both the mean and the median
          \item The median, but not the mean
          \item None of them: only one time is wrong
          \item Only the IQR
        \end{choices}
        \end{multicols}

  \boxguard[8]
  \item For 40 bottles of water labeled 500 mL, the standard deviation of the amounts is 3 mL.
        Which is the best interpretation?
        \begin{choices}
          \item Every bottle holds within 3 mL of the mean amount.
          \item A typical bottle holds 3 mL.
          \item The fullest bottle holds 3 mL more than the least full bottle.
          \item \ans{The amounts typically vary by about 3 mL from the mean amount.}
          \item About 3\% of the bottles hold more than the mean amount.
        \end{choices}

  \boxguard[8]
  \item At a checkup, a baby's length is at the \textbf{30th percentile} for babies of the same
        age. This means
        \begin{choices}
          \item the baby is 30 centimeters long.
          \item the baby is 30\% shorter than a typical baby of the same age.
          \item about 70\% of babies the same age are this length or shorter.
          \item \ans{about 30\% of babies the same age are this length or shorter.}
          \item the baby is longer than exactly 30 other babies.
        \end{choices}

  \boxguard[9]
  \item \begin{minipage}[t]{0.56\linewidth}\raggedright
        The boxplot shows the number of skaters at a skate park on each of 40 days. Which
        statement is true?
        \end{minipage}\hfill
        \begin{minipage}[t]{0.40\linewidth}
        \centering
        \begin{tikzpicture}[x=0.15cm, y=1cm, baseline=(current bounding box.north)]
          \draw[thick] (10,0) -- (50,0);
          \foreach \x in {10,20,...,50} {
            \draw (\x,-0.07) -- (\x,0.07); \node[below, font=\tiny] at (\x,-0.05) {\x};}
          \draw (12,0.45) -- (20,0.45);
          \draw (36,0.45) -- (48,0.45);
          \draw[fill=skymid] (20,0.25) rectangle (36,0.65);
          \draw[thick] (31,0.25) -- (31,0.65);
          \draw (12,0.33) -- (12,0.57);
          \draw (48,0.33) -- (48,0.57);
          \node[font=\tiny] at (30,-0.45) {Skaters per day};
        \end{tikzpicture}
        \end{minipage}
        \begin{choices}
          \item More days had 20 to 31 skaters than 31 to 36, because that part of the box is wider.
          \item About half of the days had more than 36 skaters.
          \item The mean number of skaters per day was 31.
          \item The range of the counts is 16 skaters.
          \item \ans{About a quarter of the days had 20 to 31, and about a quarter 31 to 36.}
        \end{choices}

  \boxguard[10]
  \item \begin{minipage}[t]{0.50\linewidth}\raggedright
        The parallel boxplots show the heights (cm) of tomato plants grown with two fertilizers.
        Which comparison is best?
        \end{minipage}\hfill
        \begin{minipage}[t]{0.46\linewidth}
        \centering
        \begin{tikzpicture}[x=0.15cm, y=1cm, baseline=(current bounding box.north)]
          \draw[thick] (25,0) -- (65,0);
          \foreach \x in {30,40,...,60} {
            \draw (\x,-0.07) -- (\x,0.07); \node[below, font=\tiny] at (\x,-0.05) {\x};}
          \draw (30,1.0) -- (38,1.0); \draw (52,1.0) -- (60,1.0);
          \draw[fill=skymid] (38,0.82) rectangle (52,1.18); \draw[thick] (45,0.82) -- (45,1.18);
          \draw (30,0.9) -- (30,1.1); \draw (60,0.9) -- (60,1.1);
          \node[left, font=\tiny] at (25,1.0) {X};
          \draw (35,0.45) -- (47,0.45); \draw (57,0.45) -- (62,0.45);
          \draw[fill=goldbg] (47,0.27) rectangle (57,0.63); \draw[thick] (52,0.27) -- (52,0.63);
          \draw (35,0.35) -- (35,0.55); \draw (62,0.35) -- (62,0.55);
          \node[left, font=\tiny] at (25,0.45) {Y};
          \node[font=\tiny] at (45,-0.45) {Plant height (cm)};
        \end{tikzpicture}
        \end{minipage}
        \begin{choices}
          \item Fertilizer X has the higher median, because its box starts farther left.
          \item The fertilizers are equally variable, because the boxes are the same height.
          \item \ans{Y has the higher median (52 vs.\ 45 cm) and the smaller IQR (10 vs.\ 14).}
          \item Fertilizer X's plants are taller, because its minimum is smaller.
          \item Every plant given Fertilizer Y is taller than every plant given Fertilizer X.
        \end{choices}

  \boxguard[8]
  \item A university reports the mean GPA of all 12,000 of its students is 3.12. A club surveys
        40 students; their mean GPA is 3.35. Which is correct?
        \begin{multicols}{2}
        \begin{choices}
          \item 3.12 is $\bar x$; 3.35 is $\mu$
          \item Both are parameters
          \item \ans{3.12 is $\mu$; 3.35 is $\bar x$}
          \item Both are statistics, since both are means
          \item 3.35 is $\mu$, since it is larger
        \end{choices}
        \end{multicols}

  \boxguard[9]
  \item Leo scored 62 on a math contest (mean 50, standard deviation 8). Kim scored 45 on a
        spelling contest (mean 33, standard deviation 10). Who did better \emph{relative to
        their own contest}?
        \begin{multicols}{2}
        \begin{choices}
          \item Equal: both are 12 points above
          \item \ans{Leo: $z = 1.5$ vs.\ Kim's 1.2}
          \item Kim: $z = 1.5$ vs.\ Leo's 1.2
          \item Leo: 62 is greater than 45
          \item Can't compare different subjects
        \end{choices}
        \end{multicols}

\end{enumerate}

\vspace{0.08in}

\boxguard[12]
\begin{headlinebox}{goldbg}
\small
\textbf{Section II --- Free Response} \hfill \textbf{40 points, 8 each}\\
Show your reasoning. Answer \emph{in context} --- name the variables and their units.
\end{headlinebox}

\vspace{0.06in}

\begin{enumerate}[leftmargin=1.9em, itemsep=0.10in, resume]

  \boxguard[8]
  \item \textbf{(8 points)} A marina keeps these records on the boats docked there.

        \vspace{3pt}
        {\centering\small
        \renewcommand{\arraystretch}{1.05}
        \begin{tabular}{@{} c l c c l @{}}
        \textbf{Slip \#} & \textbf{Boat type} & \textbf{Length (ft)} & \textbf{Passengers allowed} &
        \textbf{Hull color} \\ \hline
        12 & Sailboat  & 28.5 & 6  & White \\
        15 & Motorboat & 22.0 & 8  & Blue \\
        21 & Sailboat  & 34.0 & 8  & White \\
        26 & Pontoon   & 24.5 & 12 & Gray \\
        \end{tabular}\par}

        \vspace{3pt}
        \textbf{(a)} What are the individuals? Is \emph{Sailboat} a variable or a value?\\[2pt]
        \ansline{The boats docked at the marina (one per row). \emph{Sailboat} is a value of \emph{Boat type}.}

        \vspace{3pt}
        \textbf{(b)} Classify \emph{Slip \#}, \emph{Length}, and \emph{Passengers allowed} as
        categorical or quantitative. For each quantitative one, say discrete or continuous.\\[2pt]
        \ansline{Slip \#: categorical (a label). Length: quantitative, continuous.}\\
        \ansline{Passengers allowed: quantitative, discrete (a count).}

        \vspace{3pt}
        \textbf{(c)} Is ``How long is the boat in slip 21?'' a statistical question? Explain, then
        write one that is.\\[2pt]
        \ansline{No: one answer (34.0 ft), read from one row --- no variability. Statistical:}\\
        \ansline{``How much do the lengths of the boats at the marina vary?''}

  \boxguard[14]
  \item \textbf{(8 points)} Two gyms asked their members to name their favorite group class.

        \vspace{3pt}
        \begin{minipage}[c]{0.46\linewidth}
        \centering\small
        \setlength{\tabcolsep}{4pt}\renewcommand{\arraystretch}{1.05}
        \begin{tabular}{@{} l | c c @{}}
         & \textbf{East Gym} & \textbf{West Gym} \\ \hline
        Yoga           & 105 & 72 \\
        Spin           & 90  & 36 \\
        Weights        & 75  & 54 \\
        Dance          & 30  & 18 \\ \hline
        \textbf{Total} & 300 & 180 \\
        \end{tabular}
        \end{minipage}\hfill
        \begin{minipage}[c]{0.50\linewidth}
        \centering
        \begin{tikzpicture}[x=0.62cm, y=0.0145cm]
          \foreach \y in {0,40,80,120} {
            \draw[linegray, very thin] (0,\y) -- (8.4,\y);
            \node[left, font=\tiny] at (0,\y) {\y};}
          \foreach \i/\e/\w in {0/105/72, 1/90/36, 2/75/54, 3/30/18} {
            \fill[navy]    (0.3+2*\i,0) rectangle (1.0+2*\i,\e);
            \fill[goldacc] (1.0+2*\i,0) rectangle (1.7+2*\i,\w);}
          \draw[thick] (0,0) -- (8.4,0);
          \draw[thick] (0,0) -- (0,130);
          \foreach \i/\lab in {0/Yoga, 1/Spin, 2/Weights, 3/Dance}
            \node[below, font=\tiny] at (1.0+2*\i,0) {\lab};
          \node[rotate=90, font=\tiny] at (-1.1,65) {Members};
          \fill[navy] (4.6,128) rectangle (4.9,138); \node[right, font=\tiny] at (4.9,133) {East};
          \fill[goldacc] (6.5,128) rectangle (6.8,138); \node[right, font=\tiny] at (6.8,133) {West};
        \end{tikzpicture}
        \end{minipage}

        \vspace{3pt}
        \textbf{(a)} What proportion of each gym's members chose yoga?\\[2pt]
        \ansline{East: $105/300 = 0.35$. \quad West: $72/180 = 0.40$.}

        \vspace{3pt}
        \textbf{(b)} A manager wrote: ``Yoga is more popular at East Gym --- look how much taller
        its yoga bar is.'' What is wrong with this reasoning?\\[2pt]
        \ansline{The bars are counts, and East has 300 members to West's 180. Compare}\\
        \ansline{relative frequencies: 35\% vs.\ 40\% --- West's share is larger.}

        \vspace{3pt}
        \textbf{(c)} Do \emph{most} West members prefer yoga? Do most East members prefer spin?
        Justify.\\[2pt]
        \ansline{West: no --- 40\% is less than half (most common, not a majority).}\\
        \ansline{East: no --- only $90/300 = 30\%$ chose spin.}

  \boxguard[14]
  \item \textbf{(8 points)} A farm weighed the 50 pumpkins in one row of its pumpkin patch.

        \vspace{3pt}
        \begin{minipage}[t]{0.42\linewidth}
        \centering
        \begin{tikzpicture}[x=0.165cm, y=0.15cm, baseline=(current bounding box.north)]
          \foreach \y in {0,5,10,15} {
            \draw[linegray, very thin] (2,\y) -- (38,\y);
            \node[left, font=\tiny] at (2,\y) {\y};}
          \foreach \l/\c in {4/6, 8/14, 12/12, 16/8, 20/5, 24/3, 28/0, 32/2}
            \draw[fill=skymid] (\l,0) rectangle (\l+4,\c);
          \draw[thick] (2,0) -- (38,0);
          \draw[thick] (2,0) -- (2,16);
          \foreach \x in {4,12,...,36} { \node[below, font=\tiny] at (\x,0) {\x}; }
          \node[font=\tiny] at (20,-3.6) {Weight (pounds)};
          \node[rotate=90, font=\tiny] at (-1.4,8) {Pumpkins};
        \end{tikzpicture}
        \end{minipage}\hfill
        \begin{minipage}[t]{0.55\linewidth}\raggedright
        \textbf{(a)} How many pumpkins weighed 20 pounds or more? What percent of the row is
        that?\\[2pt]
        \ansline{$5+3+0+2 = 10$ pumpkins; $10/50 = 20\%$.}

        \vspace{3pt}
        \textbf{(b)} Can the exact median be found from the histogram? If not, which interval
        holds it? Explain.\\[2pt]
        \ansline{No --- values are grouped. 20 are}\\
        \ansline{under 12 lb, 32 under 16, so the}\\
        \ansline{25th and 26th are in 12--16 lb.}
        \end{minipage}

        \vspace{3pt}
        \textbf{(c)} Describe the distribution of pumpkin weights.\\[2pt]
        \ansline{Skewed right (a long tail toward heavy pumpkins), one peak at 8--12 lb.}\\
        \ansline{A gap at 28--32 lb; the 2 pumpkins at 32--36 lb are possible outliers.}\\
        \ansline{Median in 12--16 lb (about 13 lb); weights run from about 4 to 36 lb.}

  \boxguard[14]
  \item \textbf{(8 points)} The lengths (minutes) of 15 episodes of a podcast, in order:\\[2pt]
        \hspace*{\fill}4 \ \ 22 \ \ 25 \ \ 27 \ \ 28 \ \ 30 \ \ 31 \ \ 32 \ \ 33 \ \ 35 \ \ 36
        \ \ 38 \ \ 40 \ \ 41 \ \ 44\hspace*{\fill}

        \vspace{2pt}
        \textbf{(a)} Find the five-number summary.
        \begin{work}
          \text{Min} = 4, \quad Q_1 = 27, \quad \text{Median} = 32, \quad Q_3 = 38, \quad
          \text{Max} = 44
        \end{work}

        \textbf{(b)} Use the $1.5 \times \text{IQR}$ rule to decide whether the 4-minute episode is
        an outlier. If it is, where does the lower whisker of a modified boxplot end?
        \begin{work}
          \text{IQR} = 38 - 27 = 11, \qquad Q_1 - 1.5(\text{IQR}) = 27 - 16.5 = 10.5
        \end{work}
        \ansline{4 $<$ 10.5, so the 4-minute episode is an outlier. The lower whisker ends at 22.}

        \vspace{3pt}
        \textbf{(c)} The mean length is 31.07 minutes. Which better describes a typical episode,
        the mean or the median? Explain.\\[2pt]
        \ansline{The median, 32 min: it is resistant, so the 4-minute outlier barely moves it.}\\
        \ansline{The nonresistant mean is dragged down to 31.07 by that one short episode.}

  \boxguard[14]
  \item \textbf{(8 points)} The amount of coffee a machine pours into a cup is approximately
        normal, with mean $\mu = 12$ ounces and standard deviation $\sigma = 0.4$ ounces.

        \vspace{2pt}
        \begin{minipage}[t]{0.40\linewidth}
        \centering
        \begin{tikzpicture}[x=0.62cm, y=1.5cm, baseline=(current bounding box.north)]
          \draw[thick, navy, domain=-3.6:3.6, samples=80, smooth]
            plot (\x, {exp(-\x*\x/2)});
          \draw[thick] (-3.8,0) -- (3.8,0);
          \foreach \k/\lab in {-3/10.8, -2/11.2, -1/11.6, 0/12, 1/12.4, 2/12.8, 3/13.2} {
            \draw (\k,-0.05) -- (\k,0.05); \node[below, font=\tiny] at (\k,-0.03) {\lab};}
          \node[font=\tiny] at (0,-0.38) {Amount poured (ounces)};
        \end{tikzpicture}
        \end{minipage}\hfill
        \begin{minipage}[t]{0.57\linewidth}\raggedright
        \textbf{(a)} About what percent of cups get between 11.6 and 12.4 ounces?\\[2pt]
        \ansline{$\pm 1$ SD from 12, so about 68\%.}

        \vspace{3pt}
        \textbf{(b)} About what percent get \emph{less than} 11.2 ounces?\\[2pt]
        \ansline{2 SDs below: half of 5\%, about 2.5\%.}
        \end{minipage}

        \textbf{(c)} One cup gets 12.7 ounces. Find its $z$-score and interpret it in context.
        \begin{work}
          z = \frac{12.7 - 12}{0.4} = 1.75
        \end{work}
        \ansline{This cup got 1.75 SDs more than the mean amount of 12 ounces.}

        \vspace{3pt}
        \textbf{(d)} A second machine pours with mean 16 oz and SD 0.6 oz. Which pour is more
        unusually large \emph{for its own machine}: 12.7 oz from the first or 17.2 oz from the second?
        \begin{work}
          z_{\text{first}} = \frac{12.7 - 12}{0.4} = 1.75, \qquad
          z_{\text{second}} = \frac{17.2 - 16}{0.6} = 2.0
        \end{work}
        \ansline{The 17.2-oz pour: 2.0 SDs above its mean vs.\ 1.75 for the 12.7-oz pour.}

\end{enumerate}

\end{document}
